Indonesia menghasilkan sekitar 23--48 juta ton sampah makanan per tahun \cite{bps2024lingkungan}\cite{indonesiapolicy2025}, dengan sektor UMKM kuliner sebagai salah satu kontributor karena kehilangan 10--30\% produk akibat surplus yang tidak terjual. Masalah ini menimbulkan kerugian ekonomi bagi UMKM sekaligus berkontribusi terhadap emisi gas rumah kaca dan pemborosan sumber daya. Di sisi lain, banyak konsumen membutuhkan makanan terjangkau namun ragu terhadap kualitas makanan diskon. Savora hadir sebagai \textit{food rescue marketplace} berbasis web \textit{responsive}/\textit{mobile-first} yang membantu UMKM kuliner menjual makanan surplus yang masih layak konsumsi dengan harga lebih terjangkau, sekaligus memberikan informasi kelayakan yang transparan. Platform menerapkan sejumlah fitur yang dapat didemonstrasikan secara realistis dalam konteks lomba: \textit{Food Trust Index} berbasis aturan (lima tingkat kelayakan) yang dihitung dari input UMKM, \textit{Food Score Decay} yang menurun mengikuti sisa waktu rescue menuju batas kedaluwarsa, \textit{Keyword Classification} dari ulasan yang memetakan keamanan UMKM ke badge Aman/Warning/Gawat, \textit{Smart Rescue Timer}, serta rekomendasi \textit{dynamic discount} berbasis status kelayakan. Transaksi bersifat \textit{cashless} melalui \textit{payment gateway} Midtrans (mode \textit{sandbox}) dengan \textit{service fee} 5\% sebagai model pendapatan, dilanjutkan \textit{self-pickup} dengan validasi \textit{pickup code}. Savora juga menyediakan slot iklan UMKM dan pihak ketiga, analitik dan \textit{insight} penjualan UMKM, \textit{Help Center} untuk proteksi customer, \textit{Waste Log} untuk evaluasi produksi UMKM, serta role \textit{Mitra Donasi} bagi organisasi non-profit yang diverifikasi admin. Sistem dibangun dengan Next.js pada sisi \textit{frontend} serta Go/Fiber dan PostgreSQL pada sisi \textit{backend}. Dengan pendekatan berbasis aturan yang transparan dan dapat diaudit \cite{kusumawati2025mobile}, Savora ditargetkan sebagai MVP (\textit{Minimum Viable Product}) yang memberikan dampak bagi UMKM (pendapatan tambahan dan evaluasi stok), konsumen (makanan layak dengan harga hemat dan informasi kelayakan jelas), serta lingkungan (pengurangan potensi makanan terbuang), selaras dengan target SDG 12.3 \cite{unsdg2024}\cite{fao2024state} untuk mengurangi \textit{food waste} pada tahun 2030.
